% !TeX root = ../main.tex
%%% ---------------------------------------------------------------------------
%%% 结论
%%%
%%% 最后一页正文应当是「听众只记住三句话，希望是哪三句」。
%%% 主动说出局限，比等着被问出来要好。
%%% ---------------------------------------------------------------------------
\section{结论}

\begin{frame}{结论}
  \begin{block}{三点结论}
    \begin{enumerate}
      \setlength{\itemsep}{0.6em}
      \item 特征金字塔的\alert{均匀融合}是小目标漏检的一个结构性来源
      \item 空间自适应的门控 + 可变形对齐可以同时解决权重与错位两个问题，
            且是标准金字塔的严格推广
      \item 在 \dataset{} 上 mAP 达 \qty{78.4}{\percent}，
            小目标 AP 提升 \qty{6.2}{\percent}，速度保持 \qty{31}{\fps}
    \end{enumerate}
  \end{block}

  \vspace{0.8em}
  \begin{itemize}
    \setlength{\itemsep}{0.5em}
    \item \textbf{局限}：独立性假设在层间强相关时不成立；仅在光学影像上验证
    \item \textbf{下一步}：引入显式尺度先验；在 SAR 等模态上验证迁移能力
  \end{itemize}
\end{frame}

\begin{frame}[plain, noframenumbering]
  \vfill
  \begin{center}
    {\usebeamerfont{title}\usebeamercolor[fg]{title} 谢谢\par}
    \vskip 1.2em
    {\usebeamercolor[fg]{structure}\rule{0.18\paperwidth}{0.8pt}}\par
    \vskip 1.2em
    {\small 代码与预训练模型：\url{https://github.com/example/msff-net}\par}
    \vskip 0.4em
    {\small 联系方式：\href{mailto:zhangsan@mail.ustc.edu.cn}
                          {zhangsan@mail.ustc.edu.cn}\par}
  \end{center}
  \vfill
\end{frame}
